\section*{Acknowledgments}
We thank Xu Feng, Yuan Li, Shi-Cheng Xia, Jianhui Zhang and Yong Zhao for valuable discussions.
We thank the CLS Collaboration for sharing the lattice ensembles used to perform this study.
%
The LQCD calculations were performed using the Chroma software suite~\cite{Edwards:2004sx}.
%
The numerical calculation is supported by Chinese Academy of Science CAS Strategic Priority Research Program of Chinese Academy of Sciences, Grant No. XDC01040100,   HPC Cluster of ITP-CAS, and Jiangsu Key Lab for NSLSCS.
The  setup for numerical simulations was conducted  on the $\pi$ 2.0 cluster supported by the Center for High Performance Computing at Shanghai Jiao Tong University. 
J.~Hua is supported by NSFC under grant No. 11735010 and 11947215.
%
Y.-S. Liu is supported by National Natural Science Foundation of China under grant No.11905126.
%
M.~Schlemmer  and A.~Sch\"afer were supported by the cooperative research center CRC/TRR-55 of DFG. 
%
P. Sun is supported by Natural Science Foundation of China under grant No. 11975127 as well as Jiangsu Specially Appointed Professor Program.
%
W. Wang is supported in part by Natural Science Foundation of China under grant No.   11735010, 11911530088,  by Natural Science Foundation of Shanghai under grant No. 15DZ2272100.
%
Q.-A. Zhang is supported by the China Postdoctoral Science Foundation and the National Postdoctoral Program for Innovative Talents (Grant No. BX20190207).

\section{Supplemental Materials}


\subsection{Simulation checks}\label{subsec:simchecks}

\begin{figure}[!th]
\begin{center}
\includegraphics[width=0.5\textwidth]{images/06_eff_mass.pdf}
\caption{ The dispersion relation of the pion state with the pion mass  from the 2pt function. The data up to 8$\pi/L$ ($\sim$2 GeV) can be described with the formula $E_{\pi}=\sqrt{m_{\pi}^2+c_1 P^2+c_2P^4a^2}$ with $c_1=0.9945(40)$ and $c_2=-0.0282(27)$. The deviation at 8$\pi/L$ from the continuum limit is around 2\%.
}\label{fig:dispersion}
\end{center}
\end{figure}


Fig.~\ref{fig:dispersion} shows the dispersion relation with the pion mass we used. The curve  shows the fit based on the formula $E_{\pi}=\sqrt{m_{\pi}^2+c_1 P^2+c_2P^4a^2}$, where the last term in the square root parameterizes discretization errors.   We used   momenta up to  8$\pi/L$ ($\sim$2 GeV).  The fit gives results---$c_1=0.9945(40)$ and $c_2=-0.0282(27)$--- that are consistent with the ground state energy  calculated from two point function.  It indicates only small discretization errors.Thus it is  expected that the dispersion relation can  recover the standard $E_{\pi}=\sqrt{m_{\pi}^2+P^2}$ in the continuum limit.

\begin{figure}[!th]
\begin{center}
\includegraphics[width=0.4\textwidth]{images/07a_hist_ori3pt_P3_b3_tsep8.pdf}
\includegraphics[width=0.4\textwidth]{images/07b_bin_3pt_P3_b3_tsep8.pdf}
\caption{Statistical check on the simulation, taking the form factor  with $P^z=6\pi/L$, $b_{\perp}$=3, $t_{2}$=8 and $t=t_{2}/2$ as example. The left panel shows the histogram for {$864$ configurations which are included
in the further analysis}, and the right panel shows the bin size dependence after we averaged all the measurements on each configuration.}\label{fig:statistics}
\end{center}
\end{figure}



Taking the form factor  with $P^z=6\pi/L$, $b_{\perp}$=3, $t_{2}$=8 and $t=t_{2}/2$ as example, Fig.~\ref{fig:statistics} shows the statistical check of the measurements we did. {We analysed 868 configurations and dropped 4 of them in the analysis due to very strong localized artifacts.} The left panel shows the histogram of 864 (configurations) $\times$ 48 (time {slices}) measurements. {It has been noticed that using the clover action with light mass and/or a coarse lattice on the dynamical configuration, the exceptional measurement, though very rare, can occur since the critical point  is not very stable. It turns out that  some  strongly localized artifacts were not observed in other CLS ensembles with finer lattice spacings but can happen in a few configurations of the coarse ensembles, for example the  A654 ensemble which we used in this analysis.  Since this is  a small portion of the total configurations, namely $4/868$, removing these configurations might be plausible. }

 After we average the measurements over the same configuration, we find that the autocorrelation effect is negligible, since no obvious bin size dependence of the result is observed, as shown in the right panel of Fig.~\ref{fig:statistics}.


\subsection{$\ell$ dependence of TMDWF}\label{subsec:ldep}


\begin{figure}[!th]
\begin{center}
\includegraphics[width=0.9\textwidth]{images/08_tmdwf_4fig_mod.pdf}
\caption{Figure for the $\ell$ dependence of $|\phi_\ell(0,b_\perp,P^z,\ell)|$ with $\{P^z, b_{\perp}, t\}=\{6\pi/L, 3a, 6a\}$ (top left, the case shown in Fig.~\ref{fig:L_dependence}), $\{P^z, b_{\perp}, t\}=\{8\pi/L, 3a, 6a\}$ (top right), $\{P^z, b_{\perp}, t\}=\{6\pi/L, 6a, 6a\}$ (bottom left), and $\{P^z, b_{\perp}, t\}=\{6\pi/L, 3a, 8a\}$ (bottom right).}\label{fig:L_dependence2}
\end{center}
\end{figure}

%%%%%%%%%%


In Fig.~\ref{fig:L_dependence2}, we give the $\ell$ dependence of $|\phi_\ell(0,b_\perp,P^z,\ell)|$ for a few more cases, similar to the $\{P^z, b_{\perp}, t\}=\{6\pi/L, 3a, 6a\}$ case shown in Fig.~\ref{fig:L_dependence} but with larger $P^z$, $b_{\perp}$ and also $t$.


\begin{figure}[!th]
\begin{center}
\includegraphics[width=0.45\textwidth]{images/09a_HeatMap_b0.pdf}
\includegraphics[width=0.45\textwidth]{images/09b_HeatMap_b3.pdf}
\includegraphics[width=0.45\textwidth]{images/09c_HeatMap_b6.pdf}
\includegraphics[width=0.45\textwidth]{images/09d_heatmap_of_renormTMDWF.pdf}
\caption{{Heatmap for the nonperturbative renormalization factor with $P=1.58$GeV ($b_{\perp}=0$: upper left, $b_{\perp}=3a$: upper right, $b_{\perp}=6a$: lower left) and matrix elements (lower right, exhibited by the $b_{\perp}=\{0,~3a,~6a\}$ sequence). In this figure, $\Gamma$ indicates the Lorentz structure in the operators while ${\cal P}$ denotes the projection. }}\label{fig:HeatMap_NPR}
\end{center}
\end{figure}


\subsection{Estimate of Operator Mixing for TMDWF}\label{subsec:mixing}

{
In order to estimate the operator-mixing effects,  we adopt the same method as Refs.~\cite{Shanahan:2019zcq,Shanahan:2020zxr} and calculate the nonperturbative RI/MOM renormalization/mixing factors,
\begin{align}
{\cal M}_{{\cal P}\Gamma}=\textrm{Tr}\left[{\cal P}\Big\langle q(p)\Big|\overline q_1\left(\frac{z }{2}n^z +\vec b\right)\Gamma \,{\cal W}(\vec b,\ell)q_2\left(-\frac{z}{2}n^z \right)\Big| q(p)\Big\rangle\right],
\end{align}
where $\Gamma$ indicates the Lorentz structure in the operators while ${\cal P}$ denotes the projection. The relative mixing effect is considered using the ratio ${\cal M}_{{\cal P}\Gamma}/{\cal M}_{\Gamma\Gamma}$. The results with three transverse separations $b_\perp=(0a, 3a, 6a)\hat{n}_x$ and off-shell quark momentum $p=(p_t,p_x,p_y,p_z)=(3.16,0,1.58,0)$ GeV are shown with the heatmap in Fig.~\ref{fig:HeatMap_NPR} (upper left, upper right and lower left panels correspondingly). The mixing effect grows with increasing $b_\perp$. This pattern is   consistent with the perturbative calculation in Ref.~\cite{Shanahan:2019zcq}.}

{
The relative mixing effect in the quasi-TMDWF can be estimated through the product of the bare quasi-TMDWF with given Lorentz structure ${\cal P}$ and the corresponding mixing factor ${\cal M}_{\Gamma{\cal P}}$,
\begin{align}
\frac{\delta_{{\cal P}}\phi(z,b_\perp,P^z)}{\phi(z,b_\perp,P^z)}\simeq\frac{{\cal M}_{\Gamma_{\phi}{\cal P}}\Big\langle 0\Big|\overline q_1\left(\frac{z }{2}n^z +\vec b\right){\cal P}\,{\cal W}(\vec b,\ell)q_2\left(-\frac{z}{2}n^z \right)\Big| {\pi(\vec{P})}\Big\rangle}{\textrm{Re}[{\cal M}_{\Gamma_{\phi}\Gamma_{\phi}}\Big\langle 0\Big|\overline q_1\left(\frac{z }{2}n^z +\vec b\right){\Gamma_{\phi}}\,{\cal W}(\vec b,\ell)q_2\left(-\frac{z}{2}n^z \right)\Big| {\pi(\vec{P})}\Big\rangle]},
\end{align}
with $\Gamma_{\phi}=\gamma_5\gamma_t$.
We give the results in the lower right panel of Fig.~\ref{fig:HeatMap_NPR} for $b_\perp=(0a, 3a, 6a)\hat{n}_x$. From the figure, one can find that the operator-mixing effects can reach   order 5\% for the transverse separation $b\sim$ 0.6 fm, while it is less significant for smaller transverse separations. 
}

%%%%%%%%%%%%%%%%%%%
%%%%%%%%%%%%%%%%%%%


%%%%%%%%%%%%%%%%%%%
%\begin{figure}[!th]
%\begin{center}
%\includegraphics[width=0.48\textwidth]{images/HeatMap_gammatgamma5_with_P3_b0.pdf}
%\includegraphics[width=0.48\textwidth]{images/HeatMap_gammatgamma5_with_P3_b3.pdf}
%\includegraphics[width=0.48\textwidth]{images/HeatMap_gammatgamma5_with_P3_b6.pdf}
%\caption{Heatmap for the matrix elements. (b=0, b=3, b=6)  }\label{fig:HeatMap_all}
%\end{center}
%\end{figure}
%%%%%%%%%%%%%%%%%%%
 

\subsection{Tabulated results for the intrinsic soft function}\label{subsec:tabulated}

\begin{figure}[!th]
\includegraphics[width=0.45\textwidth]{images/04_sf_perturb_sqrt2.pdf}
\includegraphics[width=0.45\textwidth]{images/14_sf_normalscale.pdf}
\caption{ The intrinsic soft factor as a function of $b_{\perp}$  with  $b_{\perp,0}=a$ as {$S_{I,\overline{\rm MS}}(b_\perp,\mu)=\frac{F(b_\perp,P^z)}{F(b_{\perp,0},P^z)}\frac{|\phi(0,b_{\perp,0},P^z)|^2}{|\phi(0,b_\perp,P^z)|^2}+{\cal O}(\alpha_s, \gamma^{-2})$}. With different pion momentum $P^z$, the results are consistent with each other. The dashed  curve shows  the result of the 1-loop calculation, {$S_{I,\overline{\rm MS}}(b_{{\perp}},\mu)=1-\frac{\alpha_sC_F}{\pi}\ln\frac{\mu^2 b_{{\perp}}^2}{4 e^{-2\gamma_E}}+{\cal O}(\alpha_s)$},  with the strong coupling constant $\alpha_s(1/b_{\perp})$. {The shaded band corresponds to the scale uncertainty  of $\alpha_s$: $\mu\in [1/\sqrt{2},\sqrt{2}]\times 1/b_\perp$. The systematic uncertainty from the operator mixing has been taken into account. {Both the panels show the same data but the left panel uses the log scale and the right panel uses the normal one.}}}\label{fig:quasi-soft-factor2}
\end{figure}
 
{For ease of comparison as given in Fig.~\ref{fig:quasi-soft-factor2}, we give a tabulated results for the intrinsic soft function in Tab.~\ref{table:soft-function}.  The perturbative results for $b\le3a$ are consistent with our calculation taking into account the errors from the scale  {dependence} in the strong coupling constant $\alpha_s(\mu)$:  $\mu\in [1/\sqrt{2},\sqrt{2}]\times 1/b_\perp$.  }


\begin{table}[http]
\begin{center} 
\caption{Tabulated results for the intrinsic soft function as shown in Fig.~\ref{fig:quasi-soft-factor}. }\label{table:soft-function}
\begin{tabular}{|c|ccccccc|cc}
\hline
 & $b=a$         & $b=2a$         & $b=3a$         & $b=4a$         & $b=5a$         & $b=6a$         & $b=7a$         \\
\hline
$P=1.05$ GeV                   & 1.000(8)  & 0.567(7)  & 0.343(6)  & 0.224(6)  & 0.153(6)  & 0.106(7)  & 0.071(7)  \\
$P=1.58$ GeV     & 1.000(20) & 0.557(17) & 0.329(13) & 0.209(14) & 0.142(18) & 0.099(22) & 0.063(25) \\
$P=2.11$ GeV & 1.000(29) & 0.571(62) & 0.374(63) & 0.223(52) & 0.119(50) & 0.043(46) & 0.047(84)  \\
\hline
%pQCD [1/2, 2] & $1.000_{-0.009}^{+0.020}$   &  $0.703_{-0.326}^{+0.093}$  &  $0.303_{-3.032}^{+0.298}$  & -  & - & - & - \\
pQCD & $1.000_{-0.008}^{+0.005}$   &  $0.703_{-0.100}^{+0.057}$  &  $0.303_{-0.461}^{+0.193}$  & $-0.334_{-2.142}^{+0.498}$  & - & - & - \\
\hline
\end{tabular}
\end{center}
\end{table} 

\subsection{Two-state fit of the form factors}\label{subsec:twostate}


In this work, we perform the following joint fit to obtain the norm of the subtracted quasi-TMDWF $|\phi_\ell(0,b_\perp,P^z,\ell)|$ and soft factor $S_I(b_\perp)$ (with $\ell=7a$),
\begin{align}
&\frac{C_3(b_\perp,P^z,t_{\rm sep},t)}{C_2(0,P^z,0,t_{\rm sep})}=\frac{|\tilde{\phi}_\ell(0,b_\perp,P^z,\ell)|^2\tilde{S}_I(b_\perp)+C_1(e^{-\Delta E t}+e^{-\Delta E (t_{\rm sep}-t)})+C_2e^{-\Delta E t_{\rm sep}}}{1+C_0 e^{-\Delta E t_{\rm sep}}}, \nonumber\\
&\frac{C_2(b_\perp,P^z,0,t)}{C_2(0,P^z,0,t)}=\frac{|\tilde{\phi}_\ell(0,b_\perp,P^z,\ell)|e^{\theta(b_\perp,P^z,\ell)}(1+C_3e^{-\Delta E t})}{1+C_0e^{-\Delta E t}},
\label{eq:fit}
\end{align}
where
\begin{align}
\tilde{\phi}_\ell(0,b_\perp,P^z,\ell)=\frac{\phi_\ell(0,b_\perp,P^z,\ell)}{\phi_\ell(0,0,P^z,0)},\ \tilde{S}_I(b_\perp)=\frac{{\phi}_L(0,0,P^z,0)A_w}{2EA_p}S_I(b_\perp),
\end{align}
and $\theta(b_\perp,P^z,\ell)$ is the phase of the  quasi-TMDWF.
%
The additional factor in the definition of $\tilde{S}_I$ will be cancelled by $\tilde{S}_I(b_\perp=a)$ {when we consider the following ratio,
\begin{align}\label{eq:S_ratio_app}
S_{I,\overline{\rm MS}}(b_\perp,\mu)= \left(\frac{S_I(b_{\perp},1/a)}{S_I(b_{\perp,0},1/a)}\right)S_{I,\overline{\rm MS}}(b_{\perp,0},\mu).
\end{align}}

%%%%%%%%%%%%%%%%%
\begin{figure}[!th]
\begin{center}
\includegraphics[width=0.45\textwidth]{images/10a_ff_P3b1.pdf}
\includegraphics[width=0.45\textwidth]{images/03_ff_P3b3.pdf}\\
\includegraphics[width=0.45\textwidth]{images/10c_ff_P3b5.pdf}
\includegraphics[width=0.45\textwidth]{images/10d_test_ff_P3b1.pdf}
\caption{  The ratios $C_3(b_\perp,P^z,t_{\rm sep},t)/C_2(0, P^z,0,t_{\rm sep})$ as   function of $t_{\rm sep}$ and $t$, with $\{P^z, b_{\perp}\}=\{6\pi/L, 1a\}$ (upper left panel), $\{P^z, b_{\perp}\}=\{6\pi/L, 3a\}$  (upper right panel) and $\{P^z, b_{\perp}\}=\{6\pi/L, 5a\}$ (lower left panel). {In the lower right panel, we have dropped the $t_{\rm sep}=6$a data for the case  with $\{P^z, b_{\perp}\}=\{6\pi/L, 1a\}$ and find that the fitted result  {is consistent with the case in the upper left panel.} }}\label{fig:joint_fit2}
\end{center}
\end{figure}
%%%%%%%%%%%%%%%%%

In Fig.~\ref{fig:joint_fit2}, we {show} the ratios $C_3(b_\perp,P^z,t_{\rm sep},t)/C_2(0, P^z,0,t_{\rm sep})$ with $P^z=6\pi/L$, $b_{\perp}=\{1a,3a,5a\}$, compared  with the two-state fit predictions (colored bands) and fitted ground state contribution (gray band). All of them show good agreement between data and fits.
{This agreement indicates that the systematic uncertainty from the fit-ranges is mild. As another estimate, we have dropped the $t_{\rm sep}=6$a data for the case  with $\{P^z, b_{\perp}\}=\{6\pi/L, 1a\}$ and the results are shown in  the lower right panel of  Fig.~\ref{fig:joint_fit2}. One can  find that the fitted result  {is   consistent  with the   case in the upper left panel within uncertainties.} }


%%%%%%%%%%%%%%%%
\begin{figure}[!th]
\begin{center}
\includegraphics[width=0.7\textwidth]{images/11_differential_summed_ratio.pdf}
\caption{ The ratios $C_3(b_\perp,P^z,t_{\rm sep},t)/C_2(0, P^z,0,t_{\rm sep})$ as   function of $t_{\rm sep}$ and $t$ compared with the differential summed  ratio $R(b_\perp,P^z,t_{\rm sep})$, with $\{P^z, b_{\perp}\}=\{6\pi/L, 3a\}$.}\label{fig:joint_fit3}
\end{center}
\end{figure}
%%%%%%%%%%%%%%%%

As another check, we also consider the differential summed  ratio
\begin{align}
R(b_\perp,P^z,t_{\rm sep})&\equiv SR(b_\perp,P^z,t_{\rm sep})-SR(b_\perp,P^z,t_{\rm sep}-1)=|\tilde{\phi}_\ell(0,b_\perp,P^z,\ell)|^2\tilde{S}_I(b_\perp)+{\cal O}(e^{-\Delta E t_{\rm sep}}),\ \nonumber\\
SR(b_\perp,P^z,t_{\rm sep})&\equiv \sum_{0<t<t_{\rm sep}} \frac{C_3(b_\perp,P^z,t_{\rm sep},t)}{C_2(0, P^z,0,t_{\rm sep})}.
\end{align}
As an example, we plot $R(b_\perp,P^z,t_{\rm sep})$ as function of $t_{\rm sep}$ in Fig.~\ref{fig:joint_fit3} for $\{P^z, b_{\perp}\}=\{6\pi/L, 3a\}$ and compare it with the standard two-state fit. We can see that the $R(b_\perp,P^z,t_{\rm sep})$ agree with the ground state contribution from the two state fit at large $t_{\rm sep}$.
 

\subsection{The possible imaginary part in extracting the Collins-Soper kernel}\label{subsec:imagpart}

\begin{figure}[!th]
\begin{center}
\includegraphics[width=0.44\textwidth]{images/12a_CS_kernel_re_for_compare.pdf}
\includegraphics[width=0.44\textwidth]{images/12b_CS_kernel_im.pdf}
\caption{ The real (left panel) and imaginary (right panel) parts of the Collins-Soper kernel when the approximation $C(xP^z,\mu)=1+{\cal O}(\alpha_s)$ is taken, based on the definition in Eq.~(\ref{eq:CSK_full}).}\label{fig:imaginal}
\end{center}
\end{figure}

%%%%%%%%%%%%%%
\begin{figure}[!th]
\begin{center}
\includegraphics[width=0.45\textwidth]{images/13_CS_kernel_compare.pdf}
\caption{The comparison on $K$ and $-|K'|=-\sqrt{K^2+K'^2_{\rm Im}}$. They are consistent with each other.}\label{fig:compare}
\end{center}
\end{figure}

The Collins-Soper kernel with the following definition
\begin{align}
&K'(b_\perp,\mu)=\frac{1}{\ln(P_1^z/P_2^z)}\ln\frac{C(xP_2^z,\mu) \Phi_{\overline{\rm MS}}(x,b_\perp,P_1^z,\mu)}{C(xP_1^z,\mu) \Phi_{\overline{\rm MS}}(x,b_\perp,P_2^z,\mu)}+{\cal O}(\gamma^{-2})=\frac{1}{\ln(P_1^z/P_2^z)}\ln\frac{\phi(0,b_\perp,P_1^z)}{\phi(0,b_\perp,P_2^z)}+{\cal O}(\alpha_s, \gamma^{-2})\, \nonumber\\
&\quad\quad\quad\ \ =\frac{1}{\ln(P_1^z/P_2^z)}\ln\left|\frac{\phi(0,b_\perp,P_1^z)}{\phi(0,b_\perp,P_2^z)}\right|+\frac{i}{\ln(P_1^z/P_2^z)}\left(\theta(0,b_\perp,P_1^z)-\theta(0,b_\perp,P_2^z)\right)+{\cal O}(\alpha_s, \gamma^{-2})\, \label{eq:CSK_full}
\end{align}
should be real, but the $P^z$ dependence of the phase $\theta(0,b_\perp,P^z)=\textrm{tan}^{-1}\frac{\phi_{\rm Im}(0,b_\perp,P^z)}{\phi_{\rm Re}(0,b_\perp,P^z)}$ can introduce an {imaginary} part of $K'$ when the approximation $C(xP^z,\mu)=1+{\cal O}(\alpha_s)$ is employed. Fig.~\ref{fig:imaginal} shows the real and imaginary parts   as   functions of $b_\perp$ with two combinations of $P_1/P_2$. The real part $K'_{\rm Re}=K$ (left panel)  corresponds to the definition used in  the main text.   The non-vanishing  imaginary part (right panel)  reflects the systematic    uncertainty due to imprecise matching.  $K$ and $-|K'|=-\sqrt{K^2+K'^2_{\rm Im}}$ are  still consistent within the statistical uncertainty of $K$ as in Fig.~\ref{fig:compare}.

To estimate the effect of inaccurate matching, we consider $|K'_{\rm Im}|$ as a systematic uncertainty and {add} it with the statistical uncertainty of $K$ in quadrature.

\begin{thebibliography}{99}
%\cite{Ellis:1991qj}
\bibitem{Ellis:1991qj} 
  R.~K.~Ellis, W.~J.~Stirling and B.~R.~Webber,
  %``QCD and collider physics,''
  Camb.\ Monogr.\ Part.\ Phys.\ Nucl.\ Phys.\ Cosmol.\  {\bf 8}, 1 (1996).
  %%CITATION = CMPCE,8,1;%%
  %730 citations counted in INSPIRE as of 08 Oct 2020


%\cite{Lin:2017snn}
\bibitem{Lin:2017snn} 
  H.~W.~Lin {\it et al.},
  %``Parton distributions and lattice QCD calculations: a community white paper,''
  Prog.\ Part.\ Nucl.\ Phys.\  {\bf 100}, 107 (2018)
  doi:10.1016/j.ppnp.2018.01.007
  [arXiv:1711.07916 [hep-ph]].
  %%CITATION = doi:10.1016/j.ppnp.2018.01.007;%%
  %125 citations counted in INSPIRE as of 08 Oct 2020


%\cite{Collins:1981uk}
\bibitem{Collins:1981uk} 
  J.~C.~Collins and D.~E.~Soper,
  %``Back-To-Back Jets in QCD,''
  Nucl.\ Phys.\ B {\bf 193}, 381 (1981)
  Erratum: [Nucl.\ Phys.\ B {\bf 213}, 545 (1983)].
  doi:10.1016/0550-3213(81)90339-4
  %%CITATION = doi:10.1016/0550-3213(81)90339-4;%%
  %1218 citations counted in INSPIRE as of 08 Oct 2020


%\cite{Collins:1984kg}
\bibitem{Collins:1984kg} 
  J.~C.~Collins, D.~E.~Soper and G.~F.~Sterman,
  %``Transverse Momentum Distribution in Drell-Yan Pair and W and Z Boson Production,''
  Nucl.\ Phys.\ B {\bf 250}, 199 (1985).
  doi:10.1016/0550-3213(85)90479-1
  %%CITATION = doi:10.1016/0550-3213(85)90479-1;%%
  %1252 citations counted in INSPIRE as of 08 Oct 2020


%\cite{Ji:2004wu}
\bibitem{Ji:2004wu} 
  X.~d.~Ji, J.~p.~Ma and F.~Yuan,
  %``QCD factorization for semi-inclusive deep-inelastic scattering at low transverse momentum,''
  Phys.\ Rev.\ D {\bf 71}, 034005 (2005)
  doi:10.1103/PhysRevD.71.034005
  [hep-ph/0404183].
  %%CITATION = doi:10.1103/PhysRevD.71.034005;%%
  %601 citations counted in INSPIRE as of 08 Oct 2020


%\cite{Ji:2004xq}
\bibitem{Ji:2004xq} 
  X.~d.~Ji, J.~P.~Ma and F.~Yuan,
  %``QCD factorization for spin-dependent cross sections in DIS and Drell-Yan processes at low transverse momentum,''
  Phys.\ Lett.\ B {\bf 597}, 299 (2004)
  doi:10.1016/j.physletb.2004.07.026
  [hep-ph/0405085].
  %%CITATION = doi:10.1016/j.physletb.2004.07.026;%%
  %359 citations counted in INSPIRE as of 08 Oct 2020


%\cite{Echevarria:2015byo}
\bibitem{Echevarria:2015byo} 
  M.~G.~Echevarria, I.~Scimemi and A.~Vladimirov,
  %``Universal transverse momentum dependent soft function at NNLO,''
  Phys.\ Rev.\ D {\bf 93}, no. 5, 054004 (2016)
  doi:10.1103/PhysRevD.93.054004
  [arXiv:1511.05590 [hep-ph]].
  %%CITATION = doi:10.1103/PhysRevD.93.054004;%%
  %76 citations counted in INSPIRE as of 08 Oct 2020


%\cite{Li:2016ctv}
\bibitem{Li:2016ctv} 
  Y.~Li and H.~X.~Zhu,
  %``Bootstrapping Rapidity Anomalous Dimensions for Transverse-Momentum Resummation,''
  Phys.\ Rev.\ Lett.\  {\bf 118}, no. 2, 022004 (2017)
  doi:10.1103/PhysRevLett.118.022004
  [arXiv:1604.01404 [hep-ph]].
  %%CITATION = doi:10.1103/PhysRevLett.118.022004;%%
  %100 citations counted in INSPIRE as of 08 Oct 2020


%\cite{Ji:2019sxk}
\bibitem{Ji:2019sxk} 
  X.~Ji, Y.~Liu and Y.~S.~Liu,
  %``TMD soft function from large-momentum effective theory,''
  Nucl.\ Phys.\ B {\bf 955}, 115054 (2020)
  doi:10.1016/j.nuclphysb.2020.115054
  [arXiv:1910.11415 [hep-ph]].
  %%CITATION = doi:10.1016/j.nuclphysb.2020.115054;%%
  %19 citations counted in INSPIRE as of 08 Oct 2020


%\cite{Ji:2013dva}
\bibitem{Ji:2013dva} 
  X.~Ji,
  %``Parton Physics on a Euclidean Lattice,''
  Phys.\ Rev.\ Lett.\  {\bf 110}, 262002 (2013)
  doi:10.1103/PhysRevLett.110.262002
  [arXiv:1305.1539 [hep-ph]].
  %%CITATION = doi:10.1103/PhysRevLett.110.262002;%%
  %368 citations counted in INSPIRE as of 08 Oct 2020


%\cite{Ji:2014gla}
\bibitem{Ji:2014gla} 
  X.~Ji,
  %``Parton Physics from Large-Momentum Effective Field Theory,''
  Sci.\ China Phys.\ Mech.\ Astron.\  {\bf 57}, 1407 (2014)
  doi:10.1007/s11433-014-5492-3
  [arXiv:1404.6680 [hep-ph]].
  %%CITATION = doi:10.1007/s11433-014-5492-3;%%
  %157 citations counted in INSPIRE as of 08 Oct 2020


%\cite{Collins:2011zzd}
\bibitem{Collins:2011zzd} 
  J.~Collins,
  %``Foundations of perturbative QCD,''
  Camb.\ Monogr.\ Part.\ Phys.\ Nucl.\ Phys.\ Cosmol.\  {\bf 32}, 1 (2011).
  %%CITATION = CMPCE,32,1;%%
  %450 citations counted in INSPIRE as of 08 Oct 2020


%\cite{Bauer:2000yr}
\bibitem{Bauer:2000yr} 
  C.~W.~Bauer, S.~Fleming, D.~Pirjol and I.~W.~Stewart,
  %``An Effective field theory for collinear and soft gluons: Heavy to light decays,''
  Phys.\ Rev.\ D {\bf 63}, 114020 (2001)
  doi:10.1103/PhysRevD.63.114020
  [hep-ph/0011336].
  %%CITATION = doi:10.1103/PhysRevD.63.114020;%%
  %1421 citations counted in INSPIRE as of 08 Oct 2020


%\cite{Bauer:2001ct}
\bibitem{Bauer:2001ct} 
  C.~W.~Bauer and I.~W.~Stewart,
  %``Invariant operators in collinear effective theory,''
  Phys.\ Lett.\ B {\bf 516}, 134 (2001)
  doi:10.1016/S0370-2693(01)00902-9
  [hep-ph/0107001].
  %%CITATION = doi:10.1016/S0370-2693(01)00902-9;%%
  %735 citations counted in INSPIRE as of 08 Oct 2020


%\cite{Bauer:2001yt}
\bibitem{Bauer:2001yt} 
  C.~W.~Bauer, D.~Pirjol and I.~W.~Stewart,
  %``Soft collinear factorization in effective field theory,''
  Phys.\ Rev.\ D {\bf 65}, 054022 (2002)
  doi:10.1103/PhysRevD.65.054022
  [hep-ph/0109045].
  %%CITATION = doi:10.1103/PhysRevD.65.054022;%%
  %1228 citations counted in INSPIRE as of 08 Oct 2020


%\cite{Ji:2020ect}
\bibitem{Ji:2020ect} 
  X.~Ji, Y.~S.~Liu, Y.~Liu, J.~H.~Zhang and Y.~Zhao,
  %``Large-Momentum Effective Theory,''
  arXiv:2004.03543 [hep-ph].
  %%CITATION = ARXIV:2004.03543;%%
  %32 citations counted in INSPIRE as of 08 Oct 2020


%\cite{Cichy:2018mum}
\bibitem{Cichy:2018mum} 
  K.~Cichy and M.~Constantinou,
  %``A guide to light-cone PDFs from Lattice QCD: an overview of approaches, techniques and results,''
  Adv.\ High Energy Phys.\  {\bf 2019}, 3036904 (2019)
  doi:10.1155/2019/3036904
  [arXiv:1811.07248 [hep-lat]].
  %%CITATION = doi:10.1155/2019/3036904;%%
  %66 citations counted in INSPIRE as of 08 Oct 2020


%\cite{Ji:2014hxa}
\bibitem{Ji:2014hxa} 
  X.~Ji, P.~Sun, X.~Xiong and F.~Yuan,
  %``Soft factor subtraction and transverse momentum dependent parton distributions on the lattice,''
  Phys.\ Rev.\ D {\bf 91}, 074009 (2015)
  doi:10.1103/PhysRevD.91.074009
  [arXiv:1405.7640 [hep-ph]].
  %%CITATION = doi:10.1103/PhysRevD.91.074009;%%
  %67 citations counted in INSPIRE as of 08 Oct 2020


%\cite{Ji:2018hvs}
\bibitem{Ji:2018hvs} 
  X.~Ji, L.~C.~Jin, F.~Yuan, J.~H.~Zhang and Y.~Zhao,
  %``Transverse momentum dependent parton quasidistributions,''
  Phys.\ Rev.\ D {\bf 99}, no. 11, 114006 (2019)
  doi:10.1103/PhysRevD.99.114006
  [arXiv:1801.05930 [hep-ph]].
  %%CITATION = doi:10.1103/PhysRevD.99.114006;%%
  %41 citations counted in INSPIRE as of 08 Oct 2020


%\cite{Ebert:2018gzl}
\bibitem{Ebert:2018gzl} 
  M.~A.~Ebert, I.~W.~Stewart and Y.~Zhao,
  %``Determining the Nonperturbative Collins-Soper Kernel From Lattice QCD,''
  Phys.\ Rev.\ D {\bf 99}, no. 3, 034505 (2019)
  doi:10.1103/PhysRevD.99.034505
  [arXiv:1811.00026 [hep-ph]].
  %%CITATION = doi:10.1103/PhysRevD.99.034505;%%
  %40 citations counted in INSPIRE as of 08 Oct 2020


%\cite{Ebert:2019okf}
\bibitem{Ebert:2019okf} 
  M.~A.~Ebert, I.~W.~Stewart and Y.~Zhao,
  %``Towards Quasi-Transverse Momentum Dependent PDFs Computable on the Lattice,''
  JHEP {\bf 1909}, 037 (2019)
  doi:10.1007/JHEP09(2019)037
  [arXiv:1901.03685 [hep-ph]].
  %%CITATION = doi:10.1007/JHEP09(2019)037;%%
  %33 citations counted in INSPIRE as of 08 Oct 2020


%\cite{Ji:2019ewn}
\bibitem{Ji:2019ewn} 
  X.~Ji, Y.~Liu and Y.~S.~Liu,
  %``Transverse-Momentum-Dependent PDFs from Large-Momentum Effective Theory,''
  arXiv:1911.03840 [hep-ph].
  %%CITATION = ARXIV:1911.03840;%%
  %20 citations counted in INSPIRE as of 08 Oct 2020


%\cite{Vladimirov:2020ofp}
\bibitem{Vladimirov:2020ofp} 
  A.~A.~Vladimirov and A.~Schäfer,
  %``Transverse momentum dependent factorization for lattice observables,''
  Phys.\ Rev.\ D {\bf 101}, no. 7, 074517 (2020)
  doi:10.1103/PhysRevD.101.074517
  [arXiv:2002.07527 [hep-ph]].
  %%CITATION = doi:10.1103/PhysRevD.101.074517;%%
  %15 citations counted in INSPIRE as of 08 Oct 2020


%\cite{Ebert:2019tvc}
\bibitem{Ebert:2019tvc} 
  M.~A.~Ebert, I.~W.~Stewart and Y.~Zhao,
  %``Renormalization and Matching for the Collins-Soper Kernel from Lattice QCD,''
  JHEP {\bf 2003}, 099 (2020)
  doi:10.1007/JHEP03(2020)099
  [arXiv:1910.08569 [hep-ph]].
  %%CITATION = doi:10.1007/JHEP03(2020)099;%%
  %23 citations counted in INSPIRE as of 08 Oct 2020


%\cite{Shanahan:2020zxr}
\bibitem{Shanahan:2020zxr} 
  P.~Shanahan, M.~Wagman and Y.~Zhao,
  %``Collins-Soper kernel for TMD evolution from lattice QCD,''
  Phys.\ Rev.\ D {\bf 102}, no. 1, 014511 (2020)
  doi:10.1103/PhysRevD.102.014511
  [arXiv:2003.06063 [hep-lat]].
  %%CITATION = doi:10.1103/PhysRevD.102.014511;%%
  %15 citations counted in INSPIRE as of 08 Oct 2020


%\cite{Bruno:2014jqa}
\bibitem{Bruno:2014jqa} 
  M.~Bruno {\it et al.},
  %``Simulation of QCD with N$_{f} =$ 2 $+$ 1 flavors of non-perturbatively improved Wilson fermions,''
  JHEP {\bf 1502}, 043 (2015)
  doi:10.1007/JHEP02(2015)043
  [arXiv:1411.3982 [hep-lat]].
  %%CITATION = doi:10.1007/JHEP02(2015)043;%%
  %163 citations counted in INSPIRE as of 08 Oct 2020


%\cite{Shanahan:2019zcq}
\bibitem{Shanahan:2019zcq} 
  P.~Shanahan, M.~L.~Wagman and Y.~Zhao,
  %``Nonperturbative renormalization of staple-shaped Wilson line operators in lattice QCD,''
  Phys.\ Rev.\ D {\bf 101}, no. 7, 074505 (2020)
  doi:10.1103/PhysRevD.101.074505
  [arXiv:1911.00800 [hep-lat]].
  %%CITATION = doi:10.1103/PhysRevD.101.074505;%%
  %16 citations counted in INSPIRE as of 08 Oct 2020


\bibitem{supplemental}
Supplemental materials.

%\cite{Edwards:2004sx}
\bibitem{Edwards:2004sx} 
  R.~G.~Edwards {\it et al.} [SciDAC and LHPC and UKQCD Collaborations],
  %``The Chroma software system for lattice QCD,''
  Nucl.\ Phys.\ Proc.\ Suppl.\  {\bf 140}, 832 (2005)
  doi:10.1016/j.nuclphysbps.2004.11.254
  [hep-lat/0409003].
  %%CITATION = doi:10.1016/j.nuclphysbps.2004.11.254;%%
  %738 citations counted in INSPIRE as of 08 Oct 2020

\end{thebibliography}
